\section{投资委员会摘要}

半导体PCB链条仍处于结构性景气中，但板块已从主题发现阶段进入业绩验证阶段。当前策略不是全行业追高，而是围绕深南电路、沪电股份、生益科技做核心跟踪，并用Q2/中报触发器决定是否加权；胜宏科技、华正新材和鹏鼎控股等弹性资产需要等待客户链、订单和利润质量进一步验证。

\begin{houseviewbox}[本机构结论]
AI服务器、ASIC、Rubin、高速交换机和光模块继续推动PCB/CCL价值池上移，但投资机会已经从“找方向”转向“找兑现质量”。当前应优先跟踪已有财务和客户验证的核心公司，避免为尚未披露客户收入、平台份额或项目经济性的弹性资产支付过高溢价。下一次上调或下调判断的关键窗口是2026Q2和中报：收入、毛利率、订单和现金流必须同步验证。
\end{houseviewbox}

\begin{dashboardbox}[首页决策结论]
\begin{tabularx}{\textwidth}{L{2.8cm}X L{2.6cm}X}
\toprule
\textbf{行业观点} & 结构性看好，但不追高 & \textbf{组合动作} & 核心跟踪 + 触发器加权 \\
\textbf{首选方向} & 高端AI PCB、PCB+载板、高速CCL & \textbf{最大风险} & 若2026E/2027E预测下修则现价透支2028E高预测 \\
\textbf{上修条件} & Q2/中报收入、毛利、现金流同步改善 & \textbf{下修条件} & 订单弱化、毛利回落、客户链证据落空 \\
\bottomrule
\end{tabularx}
\end{dashboardbox}

这个判断的关键不是“AI PCB是否有需求”，而是“谁能把需求转成高质量利润”。按标的看：(1) 深南电路更接近确定性资产，体现PCB+载板交付证据；(2) 沪电股份体现高端数通PCB收入线索，但估值较高；(3) 生益科技是材料主线，但估值已要求M8/M9价格和高端材料占比继续兑现；(4) 胜宏科技弹性最高，但客户集中度、A/H估值可比性和预测分歧也最大；(5) 华正新材属于国产材料期权，小利润基数下既有弹性也有估值脆弱性。

\section{标的排序与动作}

排序采用四个定性权重（非可加总量化打分模型）：交付确定性40\%、客户链证据25\%、估值安全边际20\%、Q2催化/风险15\%。权重为判断性配置，目的在于声明各维度相对重要性，而非产出可逐格复核的评分卡；逐标的逐维打分底稿未在公开资料范围内建立。深南电路排第一，是因为PCB+载板的交付证据和盈利预测相对均衡；沪电股份排第二，是因为高端数通PCB确定性强但估值较高；生益科技排第三，是因为材料逻辑强但价格周期风险更突出；胜宏科技、华正新材和鹏鼎控股则分别对应高弹性客户链、国产材料期权和主题储备，排序折价来自客户/平台收入尚未完全验证。

\begin{exhibitbox}[图表1-1：首页估值锚和投资动作]
\scriptsize
\begin{tabularx}{\textwidth}{L{1.55cm}C{0.65cm}L{1.25cm}R{1.35cm}R{1.95cm}R{1.45cm}R{1.55cm}X}
\toprule
\textbf{公司} & \textbf{排序} & \textbf{动作} & \textbf{当前价} & \textbf{一致预期目标区间}\textsuperscript{1} & \textbf{相对区间空间} & \textbf{2026Q1已披露净利}\textsuperscript{2} & \textbf{投资含义} \\
\midrule
深南电路 & 1 & 核心跟踪 & 444.27元 & 288--373.12元 & -35\%/-16\% & 8.50亿元 & 验证后加权。 \\
沪电股份 & 2 & 核心跟踪 & 146.55元 & 101--115.03元 & -31\%/-22\% & 12.42亿元 & 上沿不追。 \\
生益科技 & 3 & 核心跟踪 & 180.15元 & 83--103.5元 & -54\%/-43\% & 11.58亿元 & 等材料验证。 \\
胜宏科技 & 4 & 进攻观察 & 361.70元 & 360--417元 & -0.5\%/+15\% & 12.88亿元 & 小仓位触发。 \\
华正新材 & 5 & 进攻/高风险 & 226.67元 & 本地公开源EPS区间（2026E 3.65--4.04元） & N/A & 0.31亿元 & 等回调至区间内再触发小仓位。 \\
鹏鼎控股 & 6 & 主题储备 & 119.01元 & 60--105.91元 & -50\%/-11\% & 4.63亿元 & 暂不进核心。 \\
\bottomrule
\end{tabularx}
\par{\scriptsize\color{mutedgrey}\textit{注：1. 一致预期目标区间为 Eastmoney 公开一致预期（\texttt{RPT\_WEB\_RESPREDICT}），非本机构目标价，亦非 house-view 推导；华正因 Eastmoney 未提供目标区间，改用本地浙商双源 EPS 区间（2026E 3.65--4.04 / 2027E 5.12--5.65 / 2028E 6.72元，来源 huazheng\_forecast\_refresh\_20260616.md），与其它5标的的 Eastmoney 目标价区间口径不同，不可直接横比；除胜宏外，多数核心标的当前价高于一致预期下沿、普遍低于上沿但仍处历史估值高位。2. 2026Q1 净利为公司一季报归母净利（official Q1 filing），非季度预测；季度拆分不公开，故不列 Q2E 单季预测。}}
\end{exhibitbox}

组合执行上，核心跟踪标的进入重点复核池；若估值回落且Q2业绩进入预期区间上沿，同时毛利率和现金流同步验证，可提高配置讨论优先级。进攻观察标的需要客户订单、产品认证或利润弹性出现新增证据后再上调跟踪权重；主题储备标的保留产业链跟踪，暂不进入核心推荐序列。若订单弱化、毛利回落、CAPEX后利用率下滑，或关键客户链证据仍停留在传闻层面，应下调目标估值和风险偏好。

除上述6只核心标的外，本报告全景覆盖半导体PCB链条共20只标的，按证据质量与估值确定性分为五层。分层口径与ch08图表8-1、图表8-25完全一致，便于读者从首页总览快速跳转至对应章节的深度分析。核心覆盖层6家内部进一步划分为三档：核心跟踪（深南、沪电、生益，3家）、进攻观察（胜宏、华正，2家）、主题储备（鹏鼎，1家），"核心跟踪"仅指证据密度最高的第一档，不可与"核心覆盖层"混淆。% FIX: 明确核心覆盖层与核心跟踪的层级关系，解决广义/狭义术语混淆

\begin{exhibitbox}[图表1-2：全景覆盖总览（20只标的分层）]
\scriptsize
\begin{tabularx}{\textwidth}{L{1.3cm}L{1.5cm}L{2.1cm}X}
\toprule
\textbf{代码} & \textbf{名称} & \textbf{分层} & \textbf{定位/核心逻辑} \\
\midrule
\multicolumn{4}{l}{\textit{核心覆盖层（6家）——本机构重点跟踪，享有目标价区间与动作指引}} \\ % FIX: 大组名从"核心跟踪"改为"核心覆盖层"，避免与组内"核心跟踪"子分类同名
\midrule
002463 & 沪电股份 & 核心跟踪 & 高端AI PCB/高速交换机核心标的，泰国数通产能+海外客户认证驱动。 \\
002916 & 深南电路 & 核心跟踪 & PCB+IC载板双轮驱动，无锡AI算力PCB募投+FC-BGA良率爬坡。 \\
600183 & 生益科技 & 核心跟踪 & 高速CCL材料主线，M8/M9价格与高端材料占比决定利润弹性。 \\
300476 & 胜宏科技 & 进攻观察 & HDI/高多层AI板弹性最大，ASIC客户+1.6T光模块mSAP订单驱动。 \\
603186 & 华正新材 & 进攻/高风险 & 国产材料期权，高等级CCL+CBF/BT小批量，利润基数低弹性大。 \\
002938 & 鹏鼎控股 & 主题储备 & AI服务器/光模块用板储备，臻鼎服务器/光模块收入验证后可升级。 \\
\midrule
\multicolumn{4}{l}{\textit{Core Candidate（3家）——可投资，证据边界要求更宽PE带}} \\
\midrule
002384 & 东山精密 & Core Candidate & 索尔思并表+1.6T光模块垂直一体化，Google ASIC客户链待半年报确认。 \\
603228 & 景旺电子 & Core Candidate & 1.6T光模块PCB批量出货，泰国产能投产修复利润是Q2关键。 \\
688630 & 芯碁微装 & Core Candidate & PCB直写光刻设备全球前三，H股挂牌+半年报OCF延续性决定估值。 \\
\midrule
\multicolumn{4}{l}{\textit{Watchlist（9家）——数据偏薄/PE极端/AI敞口弱，仅PE与预警跟踪}} \\ % FIX: Watchlist从8家改为9家，纳入建滔积层板
\midrule
301217 & 铜冠铜箔 & Watchlist & HVLP5+IC封装载体铜箔国产替代故事期，毛利率仅8.79\%。 \\
601208 & 东材科技 & Watchlist & AI高速树脂+Rubin预期，扣非弱于表观+董事长减持信号背离。 \\
603002 & 宏昌电子 & Watchlist & PE>4000x失效，高速CCL认证转批量节奏与成本压力需观察。 \\
603256 & 宏和科技 & Watchlist & Low-Dk电子布量价齐升，TTM PE约750x叙事溢价极值。 \\
002080 & 中材科技 & Watchlist & "AI电子布龙头"叙事但实际占比仅2.08\%，OCF与净利严重背离。 \\
600176 & 中国巨石 & Watchlist & 普通电子布/粗纱周期涨价驱动，公司澄清AI敞口为0。 \\
002636 & 金安国纪 & Watchlist & 典型伪AI概念，公司两次澄清无英伟达/华为/M6+敞口。 \\
300576 & 容大感光 & Watchlist & AI光刻胶预期高但增收不增利，深交所监管函为治理风险点。 \\
01888 & 建滔积层板 & Watchlist / HK Anchor & 全球CCL量价风向标，港股通南向买入为A股CCL/PCB估值预警信号；兼具Watchlist跟踪与HK Anchor双重定位。 \\ % FIX: 建滔从Demand Indicator组移至Watchlist组，保留HK Anchor特殊身份，与ch08图表8-25一致
\midrule
\multicolumn{4}{l}{\textit{Demand Indicator（2家）——需求侧验证与景气锚，非核心投资标的}} \\ % FIX: 从3家改为2家，建滔已移至Watchlist组
\midrule
601138 & 工业富联 & Demand Indicator & 800G/1.6T交换机+GB300放量，需求侧拉货验证锚，非持仓标的。 \\
300308 & 中际旭创 & Demand Indicator & 1.6T光模块龙头，需求侧拉货验证锚，非持仓标的。 \\
\bottomrule
\end{tabularx}
\vspace{0.3em}
\scriptsize 注：分层口径与ch08图表8-1、图表8-25一致；核心覆盖层6家享有AStock目标价区间（见图表1-1），其中核心跟踪3家为证据密度最高档；Core Candidate及以下不享有目标价区间，仅作结构化定位与预警跟踪。Demand Indicator与HK Anchor为产业链参照锚，非本机构A股持仓标的。 % FIX: "核心跟踪6家"改为"核心覆盖层6家"，明确核心跟踪仅为第一档3家
\end{exhibitbox}

\section{为什么不是全行业买入}

当前板块最大矛盾是“产业趋势强”与“估值已贵”同时存在。AI服务器确实提高高层板、HDI、背板/中板和高速材料需求。按反向估值矩阵（reverse\_valuation\_requirement\_matrix\_20260618），全部6个标的当前市值对公开高预测均呈负gap（沪电--61\%至--52\%、胜宏--59\%至--53\%、深南--45\%至--33\%、生益--29\%至--6\%、华正--44\%至--25\%、鹏鼎--24\%至--5\%），即尚未透支2028E（华正2028E、鹏鼎2027E）高预测；但除胜宏外，多数核心标的当前价高于一致预期下沿、普遍低于上沿但仍处历史估值高位，且要求2026E/2027E年度净利进入公开预测上沿，因此若2026E/2027E预测下修或Q2/中报验证不达上沿，现价即对2028E增长构成透支。此时，如果只根据需求景气给出全行业买入，会忽略三个约束：客户收入拆分不完整、Q2/中报尚未验证、二级市场交易拥挤。

\begin{exhibitbox}[图表1-3：不是全行业买入的三条约束]
\small
\begin{tabularx}{\textwidth}{L{3.0cm}X X}
\toprule
\textbf{约束} & \textbf{具体表现} & \textbf{投资含义} \\
\midrule
估值约束 & 多数核心标的PE/PB分位已处高位（腾讯快照PE TTM：深南83.3x、沪电65.6x、生益111.4x/PB 27.4x、胜宏76.0x、华正122.9x；沪电/深南/生益/华正三年PE分位≥99\%，胜宏86.8\%相对温和（来源：ch08图表8-11/8-12）），且除胜宏外，多数核心标的当前价高于一致预期下沿、普遍低于上沿但仍处历史估值高位，部分目标价上行空间被当前价格吸收。 & 不追高，采用一致预期区间和触发器加权。 \\
证据约束 & 具名客户收入、NVIDIA/ASIC/国产算力平台份额和客户/平台EPS桥仍缺。 & 不把传闻写成确定收入，不用伪精确目标价。 \\
兑现约束 & 高CAPEX、材料成本、折旧和客户认证会影响利润质量。 & 按上述触发器（收入、毛利率、现金流同步验证）执行Q2/中报判断。 \\
\bottomrule
\end{tabularx}
\end{exhibitbox}

综合上述约束，当前结论是结构性看好但选择性配置。任何上修都必须由财务兑现和客户链证据共同驱动，不能仅凭AI PCB景气叙事扩张风险偏好。

\section{环节评分与依据}

环节评分用于解释为什么上述排序成立。需求强、壁垒强并不自动意味着买入；如果估值也拥挤，投资动作应从“买入”转为“等待验证”。相反，需求中等但验证点明确的环节，可以在估值回落或订单增强时获得更好的风险收益比。

\begin{exhibitbox}[图表1-4：环节评分和投资含义]
\small
\begin{tabularx}{\textwidth}{L{2.4cm}C{1.3cm}C{1.3cm}X X}
\toprule
\textbf{环节} & \textbf{需求} & \textbf{壁垒} & \textbf{估值} & \textbf{评分依据和动作} \\
\midrule
高端AI PCB & 高 & 高 & 拥挤（沪电66x/胜宏76x/鹏鼎74x PE TTM） & AI服务器、交换机和光模块直接拉动高层板/HDI；沪电、胜宏、鹏鼎受益，但必须验证订单、良率和毛利率。 \\
高频高速CCL & 高 & 高 & 拥挤（生益111x/华正123x PE TTM） & M8/M9材料是信号完整性瓶颈；生益、南亚、华正受益，但需验证价格、交期和高端材料占比。 \\
PCB+IC载板 & 中高 & 高 & 分化（深南83x PE TTM） & 深南、兴森受益于载板国产化，但良率、折旧和客户认证决定利润转换。 \\
设备耗材 & 中 & 中 & 周期性（大族156x/芯碁181x PE TTM） & 大族、芯碁、鼎泰受益扩产周期，但订单滞后且CAPEX波动更大，适合储备跟踪。 \\
低端通用PCB & 低 & 低 & 分化 & 与AI价值池迁移相关性弱，不纳入AI PCB核心主线。 \\
\bottomrule
\end{tabularx}
\par{\scriptsize\color{mutedgrey}\textit{资料来源：本机构house view、source registry、claim audit；估值括注PE取自腾讯公开行情快照（2026-06-18抓取）。}}
\end{exhibitbox}

评分卡的结论是：本轮机会不是PCB全行业周期修复，而是价值池向高端板、低损耗材料和复杂工艺迁移。只有能证明技术升级转化为利润质量的公司，才值得进入核心跟踪。

\section{数据质量和使用边界}

现有资料已经覆盖核心/观察名单券商PDF、官方年报、官方IR记录和公开持仓/估值代理，但尚未覆盖完整付费一致预期数据库、具名客户收入拆分或全票实时资金流。这个缺口直接影响投资动作：当前适合使用一致预期区间和触发器框架，不适合把客户传闻或单一券商目标价转化为本机构单点目标价。

\begin{sourcequalitybox}[结论使用边界]
当前排序用于确定跟踪优先级。若后续取得公司公告、IR记录或原始模型表，证明客户链和利润质量增强，可上调标的分组；若Q2/中报只验证收入、不验证毛利率和现金流，或2026E/2027E预测下修使现价对2028E高预测转为透支，应下调目标估值和风险偏好。因此ch01给出的排序仅作跟踪优先级，下一次实质调整窗口为2026Q2业绩与中报；具体兑现触发器详见ch10。
\end{sourcequalitybox}
